\documentclass[11pt]{article}

\usepackage[margin=1in]{geometry}
\usepackage[T1]{fontenc}
\usepackage[utf8]{inputenc}
\usepackage{lmodern}
\usepackage{amsmath,amssymb}
\usepackage{booktabs}
\usepackage{array}
\usepackage{enumitem}
\usepackage{longtable}
\usepackage{xcolor}
\usepackage{hyperref}
\usepackage{listings}

\hypersetup{
    colorlinks=true,
    linkcolor=blue,
    urlcolor=blue
}

\lstset{
    basicstyle=\ttfamily\small,
    breaklines=true,
    columns=fullflexible,
    keepspaces=true,
    frame=single
}

\title{High-Level Overview of the \texttt{dev} Project}
\author{}
\date{}

\begin{document}
\maketitle

\section{Core Idea}

The project is a Python package for experimenting with \textbf{multi-agent pathfinding} (MAPF) under two different map-structure interpretations:
\begin{itemize}[leftmargin=2em]
    \item \textbf{classical mapping}: local free-space adjacencies are exposed as ordinary bidirectional movement,
    \item \textbf{cyclic mapping}: local movement is biased by a directed cyclic pattern and then postprocessed so that obstacle contact and connectivity remain manageable.
\end{itemize}

The project is therefore not mainly about inventing a new high-level MAPF solver. Instead, it studies what happens when the \emph{same underlying planning problem family} is presented through two different movement topologies. The current code now keeps that mapping comparison while allowing one global switch between vanilla CBS and ECBS for the whole experiment.

\section{What Is Being Compared}

The comparison is organized across four branches:
\begin{enumerate}[leftmargin=2em]
    \item \texttt{static\_artificial},
    \item \texttt{static\_campus\_area\_1},
    \item \texttt{dynamic\_port},
    \item \texttt{dynamic\_campus\_area\_2}.
\end{enumerate}

The branches differ in how maps are generated or loaded, how starts and goals are positioned, and whether the environment is static or time-varying. The common comparison axis is always:
\[
\text{same branch} \quad+\quad \text{same run configuration} \quad+\quad \text{different mapping}.
\]

\section{Main Design Principles}

\subsection{Paired Fairness}

For every \emph{reported} condition, classical and cyclic are compared on the same retained run configurations. A sampled configuration is retained only if both mappings produce a counted outcome. This is a very strong design choice and one of the most important invariants in the repository.

\subsection{Branch-Aware Semantics}

The four branches are not cosmetic labels. They differ in meaningful ways:
\begin{itemize}[leftmargin=2em]
    \item static vs.\ dynamic environments,
    \item artificial generation vs.\ image-based input,
    \item dispersed vs.\ clustered start placement,
    \item dispersed vs.\ clustered vs.\ single target placement,
    \item unrestricted free-cell sampling vs.\ zone-restricted campus semantics.
\end{itemize}

\subsection{Composite-Grid Representation}

Instead of storing only free cells and obstacles, the project stores explicit transition slots between vertices. This allows the mapping type itself to be represented directly in the grid rather than treated as a separate hidden rule.

\subsection{One-to-One Assignment Contract}

All current branches follow the classical one-agent--one-goal MAPF structure. The project no longer uses a shared-goal campus mode. The experiment varies \emph{where} starts and goals are placed, not the underlying assignment cardinality.

\subsection{Disappearing-Agent Interpretation}

After an agent reaches its goal, it is treated as absent from later timesteps for conflict purposes. This is true in both the static and dynamic solvers and strongly affects how congestion near targets behaves.

\section{Current Execution Shape}

At present, the code is best understood as a generalized branch runner with specialized preparation logic and one manual persistence layer:
\begin{enumerate}[leftmargin=2em]
    \item resolve a branch specification from \texttt{master\_config.py},
    \item optionally recompute branch-local raw MAPF data and save it under \texttt{dev/outputs/raw\_mapf\_files/<map\_type>/},
    \item prepare either a fresh static run context or a shared dynamic branch state,
    \item sample run configurations condition by condition,
    \item solve each retained configuration under both mappings using the currently selected global solver family,
    \item save the raw branch result payload,
    \item optionally regenerate either graphs/data outputs or Pillow visualizations from the saved raw payload.
\end{enumerate}

\section{What This Project Already Assumes}

The current project implicitly assumes the following classical MAPF conditions unless a branch-specific rule says otherwise:
\begin{itemize}[leftmargin=2em]
    \item discrete time,
    \item one-cell occupancy per agent,
    \item move-or-wait actions,
    \item offline planning with full knowledge of the dynamic loop,
    \item vertex and edge conflict detection,
    \item one unique goal per agent,
    \item dispersed sets avoid internal 8-neighbor adjacency, while clustered sets are controlled by the global \texttt{compact\_clustering} switch: compact directly adjacent clusters when \texttt{True}, and connected one-cell-gap clusters when \texttt{False}.
\end{itemize}

\end{document}
